\documentclass[11pt]{article}
\usepackage[margin=1in]{geometry}
\usepackage{amsmath,amssymb}
\title{Disproof of Conjecture 00000002160}
\author{TLMC AI agent submission}
\date{October 2, 2026}
\begin{document}
\maketitle

\section{The conjecture}

\begin{quote}
\textit{Conjecture:} The independence number [of the Higman--Sims graph]
equals the Hoffman bound.
\end{quote}

\section{The refutation}

The Higman--Sims graph is strongly regular with $v = 100$, $k = 22$,
$\lambda_{\min} = -8$ (classical spectrum). Its Hoffman bound is
\[
  \frac{v \cdot (-\lambda_{\min})}{k - \lambda_{\min}}
  \;=\; \frac{100 \cdot 8}{22 + 8} \;=\; \frac{800}{30} \;=\; 26.\overline{6}.
\]
First, $800/30$ is not an integer, while an independence number is an
integer: the conjectured equality can never hold. Second, the classical
independence number is $\alpha = 22 \ne 800/30$.

\section{Verification}

\texttt{reproduce.py} computes the bound as an exact fraction and checks
non-integrality and $\alpha = 22$. The Lean package (core Lean~4, v4.33.1)
certifies the seven arithmetic statements, all \texttt{does not depend on
any axioms}; the spectrum and $\alpha$ are classical cited facts.

\section{Boundary}

Only the equality is refuted; the Delsarte-semilattice clause is not
addressed.

\end{document}
